\documentclass[10pt,twocolumn]{article}

\usepackage[letterpaper,margin=0.82in,columnsep=0.25in]{geometry}
\usepackage{amsmath,amssymb}
\usepackage{booktabs}
\usepackage{listings}
\usepackage{xcolor}
\usepackage{url}
\usepackage[hidelinks]{hyperref}
\usepackage{microtype}
\usepackage{enumitem}

\setlist{nosep,leftmargin=*}
\setlength{\parskip}{2pt}
\setlength{\parindent}{1em}
\setlength{\textfloatsep}{7pt plus 2pt minus 2pt}
\setlength{\floatsep}{6pt plus 2pt minus 2pt}
\setlength{\intextsep}{6pt plus 2pt minus 2pt}

\definecolor{codegray}{rgb}{0.96,0.96,0.96}
\lstset{
  basicstyle=\ttfamily\footnotesize,
  backgroundcolor=\color{codegray},
  frame=single,
  framerule=0.2pt,
  columns=fullflexible,
  keepspaces=true,
  breaklines=true,
  xleftmargin=2pt,
  xrightmargin=2pt,
  aboveskip=4pt,
  belowskip=4pt
}

\newcommand{\tool}{\textsc{MiniSym}}
\newcommand{\MiniC}{\textsc{MiniC}}
\newcommand{\code}[1]{\texttt{\detokenize{#1}}}
\newcommand{\llbracket}{\mathopen{[\![}}
\newcommand{\rrbracket}{\mathclose{]\!]}}

\title{\tool: A Small Verified Symbolic Executor\\for a C-Like Language in Lean 4}
\author{Anonymous Author(s)}
\date{}

\begin{document}
\maketitle

\begin{abstract}
Symbolic execution engines are compact in concept but subtle in their treatment
of stores, branch conditions, and exceptional outcomes.  A tool that merely
enumerates plausible formulas provides little assurance that its paths cover the
behaviors of the source semantics.  We present \tool, a small executable symbolic
executor implemented and proved sound in Lean~4.  Its verified core targets a
imperative language with integer expressions, assignments, bounded
\code{while} loops, conditionals, sequencing, and assertions.  Assertions are first-class results:
the concrete and symbolic semantics distinguish normal termination from failure
and retain the store at the assertion site.  The central mechanized theorem says
that every concrete execution from a represented state is covered by a generated
symbolic result under the same input valuation, with the same result tag and an
equivalent store.  An unverified C-like parser and command-line interface make
the core directly executable on small \code{.c} files while maintaining an
explicit trust boundary.  The 1,114-line Lean artifact builds without proof holes,
contains executable regression examples, and keeps feasibility checking outside
the trusted executor.  \tool is intended as a research seed for certified
pruning, solver integration, and richer program semantics rather than as a
production C analyzer.
\end{abstract}

\section{Introduction}

Symbolic execution replaces concrete inputs with symbolic terms and accumulates
formulas describing the decisions made along each path~\cite{king1976}.  Modern
systems demonstrate that the technique can drive testing and program analysis at
substantial scale~\cite{klee,rosette}.  Their engineering challenges, however,
can obscure a basic semantic question: does the symbolic executor cover every
behavior admitted by the concrete language, including assertion failure?

This question is especially relevant for a research prototype.  Adding a parser,
solver, simplifier, and search heuristics before fixing the semantic contract can
make the trusted base difficult to identify.  \tool instead starts with a small
executable model in Lean~4~\cite{lean4}.  The concrete semantics, symbolic domain,
executor, and soundness proof are ordinary Lean definitions checked by the
kernel.  Solver interaction is deliberately absent from the theorem.

The contributions of this artifact are:

\begin{itemize}
  \item an executable concrete and symbolic semantics for an imperative core
        with fuel-bounded loops over mathematical integers;
  \item explicit, tag-sensitive modeling of normal and assertion-failure
        outcomes, including terminal failure through sequencing;
  \item a mechanized forward-soundness theorem connecting every concrete run to
        a generated symbolic path under the same valuation; and
  \item a usable, regression-tested C-like frontend and CLI whose unverified
        status is clearly separated from the verified AST-level core.
\end{itemize}

The scope is intentionally modest.  \tool does not implement ISO C, unbounded
loop reasoning, machine overflow, memory, solver-driven pruning, or a performance-oriented search
strategy.  Its purpose is to make the definitions and proof obligations of
symbolic execution small enough to inspect and extend.

\section{Tool Overview}

Figure~\ref{fig:pipeline} shows the architecture.  A handwritten lexer and
recursive-descent parser translate a restricted C-like surface language to the
\MiniC{} AST.  From that point, the concrete evaluator and symbolic executor are
defined over the same syntax.  Interpretation under a valuation connects
symbolic terms to concrete integers.  The dashed frontend boundary is important:
the main theorem begins at the AST and is not a parser-correctness theorem.

\begin{figure}[t]
\centering
\small
\setlength{\fboxsep}{4pt}
\begin{tabular}{c}
\fbox{C-like text \quad $\xrightarrow{\text{unverified parser}}$ \quad \MiniC{} AST}\\[3pt]
$\swarrow$ \hspace{48mm} $\searrow$\\[-2pt]
\fbox{\code{exec}: concrete result}
\hspace{3mm}
\fbox{\code{symExec}: result list}\\[3pt]
$\phantom{xxx}\underbrace{\hspace{57mm}}_{\text{Matches / ResultMatches under valuation }\nu}$\\[-3pt]
\multicolumn{1}{r}{\scriptsize future feasibility pruning $\dashrightarrow$}
\end{tabular}
\caption{Architecture and trust boundary.  Feasibility is downstream of the
verified executor and is not assumed by soundness.}
\label{fig:pipeline}
\end{figure}

The following accepted input illustrates the interface:

\begin{lstlisting}[language=C]
int x = input;
if (x > 0) {
  int y = x + 1;
  assert(y != 0);
} else {
  int y = x - 1;
  assert(y != 0);
}
\end{lstlisting}

Running \code{lake exe minisym examples/demo.c} parses the file, symbolically
executes it from the canonical input state, and prints selected symbolic-store
bindings with each path condition.  The example produces four raw results: each
control-flow branch reaches an assertion, and each assertion emits a normal and
a potential-failure result.  The word ``potential'' matters: without a
feasibility proof, a failure path need not correspond to an actual input.
The option \code{--fuel N} overrides the default loop bound.

The parser accepts initialized \code{int} declarations, assignments, blocks,
mandatory \code{if}/\code{else}, \code{assert}, comments, addition and
subtraction, six comparisons, negation, conjunction, disjunction, and
\code{while} loops.  It lowers
\code{!=} to negated equality and reverses operands for \code{>} and \code{>=}.
An initialized declaration is only assignment syntax; there are no types,
scopes, or shadowing rules.  Uninitialized declarations and unsupported syntax
such as \code{for} are rejected with line and column diagnostics.

\section{Executable Semantics}

\subsection{Core Language and Concrete Execution}

Arithmetic expressions are constants, variables, addition, and subtraction.
Boolean expressions contain equality, strict and non-strict order, negation,
conjunction, and disjunction.  Statements are \code{skip}, assignment,
sequencing, conditional choice, bounded \code{while}, and assertion.  Execution
is total because a natural-number fuel parameter bounds recursive loop unrolling.

A concrete store is a total function
$\sigma : \mathsf{String}\rightarrow\mathbb{Z}$.  Thus reading a variable always
returns the value supplied by the incoming store; the model has no notion of an
uninitialized read.  Functions $\llbracket e\rrbracket_\sigma$ and
$\llbracket b\rrbracket_\sigma$ evaluate expressions to an integer and a Boolean.
Execution has the type
\[
 \mathsf{exec}:\mathsf{Stmt}\rightarrow\mathsf{State}
                  \rightarrow\mathsf{ExecResult},
\]
where a result is $\mathsf{normal}(\sigma')$,
$\mathsf{assertionFailure}(\sigma_f)$, or
$\mathsf{fuelExhausted}(\sigma_e)$.  An assertion checks its Boolean
expression and retains the current store in either result.  Failure is terminal:
if the first component of a sequence fails, the second is not evaluated.  This
explicit failure store is essential for stating correspondence at the failure
point.  A false loop guard exits normally.  A true guard with positive fuel
executes one body iteration and continues with decremented fuel; a true guard at
zero fuel returns \code{fuelExhausted}.  Nested loops receive the current budget,
so the parameter is not a global bound on the sum of all iterations.

\subsection{Symbolic Domain}

Symbolic arithmetic terms form a separate AST with integer constants, named
inputs, addition, and subtraction.  Likewise, path conditions use an explicit
\code{SymBool} syntax rather than Lean propositions.  This representation can be
printed, traversed, normalized, or eventually translated to SMT-LIB without
changing the core semantics.

A valuation $\nu:\mathsf{String}\rightarrow\mathbb{Z}$ interprets named inputs.
The functions $\mathcal I_E(\nu,t)$ and $\mathcal I_B(\nu,\phi)$ interpret
symbolic terms and formulas.  A symbolic state
$S=\langle\mu,\pi\rangle$ contains a total symbolic store
$\mu:\mathsf{String}\rightarrow\mathsf{SymExpr}$ and a path condition $\pi$.
Initially,
\[
  \mu_0(x)=\mathsf{input}(x), \qquad \pi_0=\mathsf{true}.
\]
Functional stores simplify pointwise proofs.  Since functions have no useful
default textual representation, the output layer takes a finite list of variable
names and creates a readable snapshot; this layer is not involved in soundness.

\subsection{Symbolic Execution and Assertions}

The executor returns all unpruned outcomes under an explicit bound:
\begin{align*}
 \mathsf{symExecFuel}:\;&\mathbb{N}\rightarrow\mathsf{Stmt}\rightarrow\mathsf{SymState}\\[-2pt]
 &\rightarrow\mathsf{List}\;\mathsf{SymResult}.
\end{align*}
Assignment symbolically evaluates the right-hand side and updates $\mu$.
Sequencing maps its second statement over normal results but preserves a failure
as a singleton, matching concrete short-circuiting.  For a guard $b$, let
$\widehat b$ be its symbolic evaluation.  A conditional appends the results of
the two executions
\[
 (s_t,\langle\mu,\pi\land\widehat b\rangle),\qquad
 (s_f,\langle\mu,\pi\land\neg\widehat b\rangle).
\]
An assertion emits exactly two tagged states:
\[
 \begin{array}{ll}
 \mathsf{normal}             & \langle\mu,\pi\land\widehat b\rangle,\\
 \mathsf{assertionFailure}   & \langle\mu,\pi\land\neg\widehat b\rangle.
 \end{array}
\]
At zero loop fuel, the false-guard path remains normal and the true-guard path is
tagged \code{fuelExhausted}; at positive fuel, the true path executes the body
and continues with decremented fuel.  Default wrappers use a bound of eight.
The executor does not ask whether either formula is satisfiable.  This choice
keeps external solvers out of the trusted argument and permits a direct proof,
at the cost of infeasible results and exponential path growth.

\section{Mechanized Soundness}

\subsection{Representation Relations}

The central invariant says when a symbolic state denotes a concrete state under
a valuation:
\begin{align*}
&\mathsf{Matches}(\nu,\sigma,\langle\mu,\pi\rangle) \;\triangleq\\
&\quad \mathcal I_B(\nu,\pi)=\mathsf{true}
 \ \land\ \forall x.\;\sigma(x)=\mathcal I_E(\nu,\mu(x)).
\end{align*}
Both conjuncts are substantive: the valuation must satisfy the path condition,
and the interpreted symbolic store must agree pointwise with the concrete store.
The lemma \code{initial_matches} proves
$\mathsf{Matches}(\nu,\nu,S_0)$ for every valuation, making the premise of the
main theorem explicitly inhabited.

Structural induction establishes exact expression agreement:
\[
 \mathsf{Matches}(\nu,\sigma,S)\Rightarrow
 \llbracket e\rrbracket_\sigma =
 \mathcal I_E(\nu,\mathsf{symEvalExpr}(S.\mu,e)),
\]
and analogously for Boolean truth values.  These results yield three preservation
lemmas: assignment updates concrete and symbolic stores in lockstep; a concretely
true guard satisfies $\pi\land\widehat b$; and a concretely false guard satisfies
$\pi\land\neg\widehat b$.

Store agreement alone would allow an incorrect symbolic outcome tag to witness a
concrete result.  \tool therefore defines \code{ResultMatches} with only two
constructor-aligned rules for normal termination, assertion failure, and fuel
exhaustion.  Writing $M$ for \code{Matches} and $RM$ for
\code{ResultMatches}, they are:
\[
 \frac{M(\nu,\sigma,S)}
      {RM(\nu,\mathsf{normal}(\sigma),\mathsf{normal}(S))}
\]
\[
 \frac{M(\nu,\sigma,S)}
      {RM(\nu,\mathsf{failure}(\sigma),\mathsf{failure}(S))}.
\]
Consequently, a concrete assertion failure cannot be represented by a symbolic
normal result, or conversely.

\subsection{Main Theorem}

The exact Lean theorem is:

\begin{lstlisting}
theorem symExecFuel_sound (fuel : Nat)
    (statement : Stmt)
    {valuation : Valuation} {state : State}
    {symbolic : SymState}
    (h : Matches valuation state symbolic) :
    exists result,
      result `elem` symExecFuel fuel statement symbolic /\
      ResultMatches valuation
        (execFuel fuel statement state) result
\end{lstlisting}

The listing uses ASCII approximations for Lean's existential, membership, and
conjunction notation.  Formally, the theorem states
\begin{align*}
 &\mathsf{Matches}(\nu,\sigma,S)\Rightarrow\\[-2pt]
 &\quad\exists R\in\mathsf{symExecFuel}(k,s,S).\\[-2pt]
 &\qquad \mathsf{ResultMatches}(\nu,\mathsf{execFuel}(k,s,\sigma),R).
\end{align*}
For a fixed valuation and a represented input-state pair, the unique concrete
outcome is therefore covered by at least one generated symbolic result under the
same valuation.  The witness has the same normal/failure/exhaustion tag, its path condition
is satisfied, and its interpreted store equals the concrete final or
failure-point store for every variable.

The proof uses outer induction on fuel and inner induction on $s$, generalizing
both states.  The \code{skip} case
uses the unchanged state, and assignment uses update preservation.  In a
sequence, the first induction hypothesis yields a matching symbolic result.  A
normal intermediate result feeds the second induction hypothesis; list
membership follows through \code{flatMap}.  A failing intermediate result is
preserved without executing the suffix.  A conditional case-splits on concrete
guard evaluation, strengthens the matching relation using the appropriate path
lemma, and embeds the branch witness in the corresponding side of list append.
The assertion cases select the emitted normal or failure result according to the
concrete Boolean value.  A loop case considers a false exit, zero-fuel
exhaustion, and a positive-fuel iteration.  A normal body result invokes the
outer induction hypothesis at the predecessor; assertion failure or nested fuel
exhaustion is propagated with its matching tag.  The theorem
\code{symExec_sound} remains as the default-fuel specialization.

This is a forward simulation or concrete-run coverage result.  It does not claim
that every emitted path is feasible, that every potential failure is realizable,
or that a solver is complete.  These exclusions are strengths of the stated
trust boundary rather than hidden assumptions.

\section{Artifact Validation}

Table~\ref{tab:examples} summarizes executable regression programs.  The
straight-line example checks nested symbolic substitution.  A conditional emits
the guard and its negation.  The nested example deliberately retains the
contradictory condition $10<\mathit{input}\land\mathit{input}<5$.  Separately, a
Lean theorem proves the corresponding integer contradiction using \code{omega};
this is a feasibility demonstration, not an integrated pruning pass.  The
assertion example checks that two control paths become four results, exactly two
of which carry the failure tag.

\begin{table}[t]
\centering
\small
\caption{Executable semantic regressions.}
\label{tab:examples}
\begin{tabular}{@{}lccp{2.9cm}@{}}
\toprule
Scenario & Paths & Fail. & Checked behavior\\
\midrule
Straight line & 1 & 0 & symbolic assignment/substitution\\
Conditional & 2 & 0 & guard and negated guard\\
Nested branch & 3 & 0 & infeasible path retained\\
Assertion & 4 & 2 & tag and path-condition split\\
Bounded loop & 5 & 0 & four exits plus one exhaustion\\
\bottomrule
\end{tabular}
\end{table}

The parser suite contains four accepted structural cases covering arithmetic
associativity, Boolean precedence and grouping, nearest-\code{if} association,
comments, declarations, and assertions.  Five negative guards check rejection of
an uninitialized declaration, a loop, a chained comparison, trailing input, and
an unterminated comment.  These are frontend regressions, not a proof of parser
correctness.

The artifact pins \code{leanprover/lean4:v4.32.2}.  Reproduction requires:
\begin{lstlisting}[language=bash]
lake build
lake exe minisym examples/demo.c
lake exe minisym --fuel 3 examples/loop.c x
lake exe minisym examples/demo.c x y input
\end{lstlisting}
The default build checks the library, executable examples, parser guards,
soundness proofs, and CLI (18 build jobs in our environment).  The project has
1,114 lines of Lean: 683 lines for syntax, the two semantics, executor, and
soundness; 235 for the parser; and the remainder for examples, tests, and CLI.
A source scan finds no \code{sorry}, \code{axiom}, \code{admit}, or \code{unsafe}
in Lean files.  These measurements establish artifact consistency, not runtime
scalability; no performance claim is made.

\section{Discussion and Research Directions}

\paragraph{Limitations.}
The verified language omits unbounded loop semantics, procedures, multiplication, arrays, heap
memory, pointers, bounded integers, overflow, and undefined behavior.  The
surface parser is C-like rather than C: declarations have no binding semantics,
\code{else} is mandatory, and the theorem begins only after parsing.  Functional
stores require users to select names for printing.  Without simplification or
feasibility pruning, the result set grows exponentially and contains false-positive
failure paths.

\paragraph{Certified feasibility.}
The explicit \code{SymBool} AST gives a clean next step: define a pruning
interface that removes a path only when accompanied by a Lean proof that its
interpretation is false.  Lean's \code{omega} tactic can already discharge some
linear-integer contradictions.  For larger formulas, an SMT translation could
remain outside the executor, with solver answers accepted only after certificate
checking.  Thus optimization need not enlarge the semantic trusted base.

\paragraph{Richer intermediate languages.}
An SSA transformation could reduce functional-store complexity and simplify SMT
output, provided it comes with a semantics-preservation theorem.  Beyond that,
finite-width arithmetic, memory regions, aliasing, and pointer provenance are
necessary before claiming C or LLVM semantics.  The present proof decomposition
also suggests a reusable abstraction: expression interpretation, store update,
branch strengthening, and constructor-sensitive result matching can potentially
be factored once for multiple frontends.

\section{Related Work}

King's formulation established symbolic execution as program execution over
symbolic values and path constraints~\cite{king1976}.  KLEE demonstrates a mature
engineering point in the design space, operating on LLVM programs and using
constraint solving to generate high-coverage tests~\cite{klee}.  Rosette embeds
solver-aided execution in a host language and supports the construction of
domain-specific tools~\cite{rosette}.  \tool does not seek feature or performance
parity with either system.  Its contribution is a small Lean artifact in which
the executor and its concrete-to-symbolic coverage theorem can be read together.

Lean~4 combines an interactive theorem prover with an efficient functional
language and extensible tooling~\cite{lean4}, allowing the semantics, executable
engine, examples, and proofs to inhabit one project.  In contrast to a verified
compiler such as CompCert~\cite{compcert}, \tool currently verifies neither a
translation from ISO C nor a low-level target.  The comparison highlights the
next major proof obligation: a richer frontend must be related formally to the
verified core rather than merely parsed into it.

\section{Conclusion}

\tool demonstrates that a useful symbolic-execution seed can be both executable
and accompanied by a non-vacuous semantic theorem.  The development makes path
conditions explicit, treats assertion failure as a genuine outcome, preserves
the same valuation and store interpretation throughout execution, and keeps
solver trust out of the core.  Its C-like CLI supports experimentation today;
its sharply stated limitations identify proof-producing pruning, certified solver
integration, SSA, and memory semantics as concrete next steps.

\clearpage
\begin{thebibliography}{9}

\bibitem{king1976}
James C. King.
Symbolic execution and program testing.
\emph{Communications of the ACM}, 19(7):385--394, 1976.

\bibitem{klee}
Cristian Cadar, Daniel Dunbar, and Dawson Engler.
KLEE: Unassisted and automatic generation of high-coverage tests for complex
systems programs.
In \emph{Proc. 8th USENIX Symposium on Operating Systems Design and
Implementation (OSDI)}, 2008.

\bibitem{rosette}
Emina Torlak and Rastislav Bodik.
A lightweight symbolic virtual machine for solver-aided host languages.
In \emph{Proc. ACM SIGPLAN Conference on Programming Language Design and
Implementation (PLDI)}, 2014.

\bibitem{lean4}
Leonardo de Moura and Sebastian Ullrich.
The Lean 4 theorem prover and programming language.
In \emph{Proc. 28th International Conference on Automated Deduction (CADE)},
pages 625--635, 2021.

\bibitem{compcert}
Xavier Leroy.
Formal verification of a realistic compiler.
\emph{Communications of the ACM}, 52(7):107--115, 2009.

\end{thebibliography}

\end{document}
